\newtoggle{withbonus} \settoggle{withbonus}{false} % set true to show bonus points in grade table

% setting underlying course name (Grundlagenfach Mathematik, §-Unterricht, \ldots)
\newcommand{\mycourse}{Informatik}
% name of the class
\newcommand{\myclass}{Klasse 2f}
% set date
\newcommand{\mydate}{19. Juni 2024}

\title{\textbf{Theoretische Informatik: randomisierte Algorithmen}}
\author{Thomas Graf\\
\mycourse, \myclass}
\date{\mydate}
\maketitle
\thispagestyle{headandfoot}		

\begin{center}
	\fbox{\parbox{140mm}{\centering
			\begin{itemize}
				\item Dauer der Prüfung: $55$ Minuten
				%\item \textbf{Zeige in jeder Aufgabe unbedingt Deine Schritte!}
				\item erlaubte Hilfsmittel:
				\begin{itemize}
				    \item Schreibzeug
            \item \textbf{ohne Taschenrechner}
				\end{itemize}
	\end{itemize}}}
\end{center}
\vspace{10mm}

\begin{center}
	Vorname: \rule{45mm}{0.8pt}\qquad Nachname: \rule{45mm}{0.8pt}
\end{center}
\vspace{20mm}

\begin{center}
	\iftoggle{withbonus} {
		\combinedgradetable[h][questions]
	} {
		\gradetable[h][questions]
	}
\end{center}
\vspace{15mm}
\begin{center}
	Note: 
	\vspace{10mm}
	
	\rule{8mm}{1.2pt}
\end{center}
%\thispagestyle{empty}
\clearpage


\begin{questions}
% 2f
\question
\begin{parts}
  \part[1] Es sei $x := 0111001$ ein binärer String. Bestimmen Sie $\text{Nummer}(x)$. 
  \part[1] Bestimmen Sie $\text{Bin}(28)$.
  \part[1] Für einen binären String $x$ gilt 
  \begin{align*}
   \text{Nummer}(x) = 2^k - 1.
  \end{align*}
  Dabei ist $k$ eine positive natürliche Zahl. Bestimmen Sie $\abs{x}$ in Abhängigkeit von $k$.
  \part[1] Bestimmen Sie $\text{Prim}(37)$.
  \part[1] Nennen Sie (konkrete) positive natürliche Zahlen $x, y$ und $z$, sodass zwar die Relation
\begin{align*}
  x < y < z
\end{align*}
gilt, aber nicht die Relation
\begin{align*}
  \abs{\text{Bin}(x)} < \abs{\text{Bin}(y)} < \abs{\text{Bin}(z)}.
\end{align*}
\part[1] Donald Knuth möchte unser randomisierte Kommunikationsprotokoll für die beiden Datensätze $x := 001110$ und $y := 1110$ durchführen. Was ist hier das Problem? Welchen Rat würden Sie Donald geben?
\end{parts}
\droptotalpoints

\question[2] In Zürich sei der Datensatz $x := 001110$ auf dem Rechner $R_1$ abgespeichert. $R_1$ wird das Kommunikationsprotokoll mit einem Rechner $R_2$ in Boston durchführen.
Dabei sei $y$ der Datensatz, welcher auf $R_2$ abgespeichert ist. In Zürich wurde zufällig die Primzahl $p := 5$ gewählt. Geben Sie ein konkretes Beispiel eines (sinnvollen) Datensatzes $y$, sodass das Protokoll falsch entscheiden wird.

\question Gegeben seien zwei Datensätze $x$ und $y$ mit einer jeweiligen Länge von $10^{24}$ Bits. 
\begin{parts}
\part[2] Geben Sie eine vernünftige obere Schranke für den Kommunikationsaufwand des Kommunikationsprotokolls zum Abgleich von $x$ und $y$ an.
\part[1] Wie viele Bits müsste ein deterministischer Algorithmus zur Beantwortung der Frage $x\stackrel{?}{=}y$ mindestens kommunizieren?
\end{parts}
\droptotalpoints

\question[2]
Für die Fehlerwahrscheinlichkeit $r$ des Protokolls haben wir die Abschätzung
\begin{align*}
    r \leq \frac{n-1}{\text{Prim}(n^2)}\leq \frac{n-1}{n^2/\ln(n^2)}\leq\frac{\ln(n^2)}{n}
\end{align*}
gefunden. Dabei ist $n$ die Länge der jeweiligen Datensätze (in Bits). Erklären Sie, warum wir im Protokoll eine Primzahl $\leq n^2$ wählen und nicht lediglich eine Primzahl $\leq n$.

\clearpage
\question[2]
Gegeben ist die (korrekte und vollständige) Ihnen vertraute Python-Funktion \pythoninline{sieb(n)}:
\begin{lstlisting}[language=Python,caption=Sieb,numbers=none]
def sieb(n):
    primzahlen = []
    nicht_gestrichen = (n + 1) * [True]

    i = 2
    while i * i <= n:
        if nicht_gestrichen[i]:
            j = 2 * i
            while j <= n:
                nicht_gestrichen[j] = False
                j += i
        i += 1

    i = 2
    while i <= n:
        if nicht_gestrichen[i]:
            primzahlen.append(i)
        i += 1

    return primzahlen
\end{lstlisting}
Verwenden Sie diese Funktion als Hilfsfunktion um eine Python-Funktion \pythoninline{prime_gaps(n)} zu definieren, welche die ersten $n$ Primzahllücken (prime gaps) printet.
\begin{lstlisting}[language=Python,caption=prime gaps (noch zu vervollständigen),numbers=none]
def prime_gaps(n):

    .............

    .............

    .............

    .............

    .............

    .............

    .............
\end{lstlisting}


% \question[2]
% Gegeben sind zwei Funktionen $f$ und $g$, welche auf ganz $\N$ definiert sind und $g(n) \neq 0$ für alle $n\in\N$. Es gelte die Gleichung
% \begin{align*}
%   \lim_{n\to\infty}\frac{f(n)}{g(n)} = 1.
% \end{align*}
% Ist es korrekt, dass daraus die Gleichung
% \begin{align*}
%   \lim_{n\to\infty} \abs{f(n) - g(n)} = 0
% \end{align*}
% folgt? Begründen Sie die Behauptung oder finden Sie ein Gegenbeispiel.


% \question[2] Es sei $f$ eine Funktion von $\N$ nach $\N$. Sei $p>0$ eine natürliche Zahl. Welche Eigenschaft muss $f$ erfüllen, damit die Implikation
% \begin{align*}
%   (\quad f(x)\;\text{mod}\; p = f(y)\;\text{mod}\; p \quad) \Rightarrow (x = )
% \end{align*}
%
% % \question[2]
% Ein Primzahltest ist ein mathematisches Verfahren, um festzustellen, ob eine gegebene natürliche Zahl $n\geq 2$ eine Primzahl ist oder nicht.
%
% Unser Primzahltest aus dem Unterricht zeigt, dass es für einen Primzahltest sicherlich ausreicht, für alle Zahlen aus der Menge $\lrc{2, 3, \ldots, \floor{\sqrt{n}}}$ zu prüfen, ob sie Teiler von $n$ sind.

\end{questions}
